%%%%%%%%%%%%%%%%%%%%%%%%%%%%%%%%%%%%%%%%%
% Tecnológico de Costa Rica/Instructivo de Laboratorio de Instrumentación I
% LaTeX Template
% Version 3.1 (25/3/14)
%
% This template has been downloaded from:
% http://www.LaTeXTemplates.com
%
% Original author:
% Linux and Unix Users Group at Virginia Tech Wiki 
% (https://vtluug.org/wiki/Example_LaTeX_chem_lab_report)
%
% License:
% CC BY-NC-SA 3.0 (http://creativecommons.org/licenses/by-nc-sa/3.0/)
%
%%%%%%%%%%%%%%%%%%%%%%%%%%%%%%%%%%%%%%%%%

%----------------------------------------------------------------------------------------
%	PACKAGES AND DOCUMENT CONFIGURATIONS
%----------------------------------------------------------------------------------------

\documentclass[12pt,letterpaper]{report}
\makeatletter
\def\input@path{{../common/}{../guide/}{../data/}{../code/}}
\makeatother
\input{../common/preamble.tex}
\input{../common/arduinoLanguage.tex}
\usepackage{csvsimple}
\usepackage{geometry} 
\geometry{left=18mm,right=18mm,top=21mm,bottom=21mm,headheight=15pt}

\setlength\parindent{0pt} % Removes all indentation from paragraphs

\renewcommand{\labelenumi}{\alph{enumi}.} % Make numbering in the enumerate environment by letter rather than number (e.g. section 6)
\usepackage{fancyhdr}
\pagestyle{fancy}

%----------------------------------------------------------------------------------------
\lhead{Instructivo de Laboratorio de Instrumentación I}
\rhead{\begin{picture}(0,0) \put(-60,0){\includegraphics[width=20mm]{logo.png}} \end{picture}}
\newcommand{\obj}{Objetivos}
\newcommand{\mat}{Materiales y equipo}
\newcommand{\pro}{Procedimiento}
\newcommand{\capacidad}{Al finalizar este laboratorio el estudiante estará en capacidad de:}
\newcommand{\antesde}{Antes de empezar el laboratorio presente el siguiente cuestionario lleno.}
%----------------------------------------------------------------------------------------
%	DOCUMENT INFORMATION
%----------------------------------------------------------------------------------------


\addto\captionsspanish{\renewcommand{\chaptername}{Laboratorio}}
\addto\captionsspanish{\renewcommand{\tablename}{Tabla}}
\begin{document}
% ------------------------------------------------------------
\begin{titlepage}
    \begin{center}
\vspace*{1in}
\begin{figure}[htb]
\begin{center}
\includegraphics[width=11cm]{../figures/logo.png}
\end{center}
\end{figure}
\vspace*{0.4in}
\begin{Large}
Escuela de Física\\
\vspace*{0.15in}
Ingeniería Física\\
\vspace*{0.8in}
\end{Large}
\vspace*{0.2in}
\begin{Large}
\textbf{Instructivo de Laboratorio} \\
\end{Large}
\vspace*{0.3in}
\begin{large}
Instrumentación I\\
\end{large}
\vspace*{2.5in}
\begin{Large}
\textbf{\today}\\
Versión: 0.2\\
\end{Large}
\rule{80mm}{0.1mm}\\
\vspace*{0.1in}
\begin{large}
Realizado por: Mía S. Cañón Gil y Kashana Porte Castro\\
\end{large}
\end{center}
\end{titlepage}
% ------------------------------------------------------------
\tableofcontents
% ------------------------------------------------------------
% Lista de practicas autogenerada por guide/build_pdfs.sh
% Archivo autogenerado por build_pdfs.sh
% Incluye todas las practicas detectadas en guide/
\input{01_terrario}
\input{02_opticos}
% !TEX root = 00_instructivo.tex
\chapter{Transductores eléctricos y magnéticos: Hall, magnetómetro y galgas}

\section*{Objetivos}
\begin{itemize}
    \item Medir campos magnéticos con sensor Hall (ACS712) y magnetómetro QMC5883L.
    \item Detectar deformación mecánica mediante galgas piezorresistivas.
    \item Integrar los sensores en un circuito sobre PCB de una cara.
\end{itemize}

\section*{Materiales y equipo}
\begin{itemize}
    \item Sensor Hall ACS712 (para corriente/campo).
    \item Magnetómetro triple eje QMC5883L (GY-271).
    \item Galgas extensiométricas (resistencia nominal 120 $\Omega$).
    \item Placa de cobre lisa (para PCB).
    \item Láminas de acrílico y polvo ferroso.
    \item Arduino Uno.
    \item Mini protoboard (solo para pruebas).
    \item Amplificador operacional LM358 (para puente de Wheatstone).
\end{itemize}

\section*{Instrucciones}
\begin{enumerate}
    \item Conecte el ACS712: VCC a 5V, GND a GND, OUT a A0.
    \item Conecte el QMC5883L mediante I2C (A4=SDA, A5=SCL).
    \item Para las galgas: monte un puente de Wheatstone con una galga activa y tres resistencias fijas. Amplifique la salida con LM358 y conéctela a A1.
    \item Coloque polvo ferroso entre acrílicos y acerque un imán. Mida el campo con el magnetómetro (componentes X,Y,Z) y con el Hall.
    \item Aplique peso sobre la galga (placa de cobre) y registre la deformación.
    \item Diseñe y fabrique una PCB (transferencia térmica) para interconectar los sensores sin protoboard.
\end{enumerate}

\section*{Esquemático}
\begin{figure}[htbp]
\centering
\begin{circuitikz}[scale=0.7]
\draw (0,0) node[above]{Arduino};
% ACS712
\draw (1,0.5) node[anchor=east]{5V} -- (2,0.5) node[anchor=west]{VCC};
\draw (1,1) node[anchor=east]{GND} -- (2,1) node[anchor=west]{GND};
\draw (1,1.5) node[anchor=east]{A0} -- (2,1.5) node[anchor=west]{OUT};
% QMC5883L I2C
\draw (1,2.5) node[anchor=east]{A4} -- (3,2.5);
\draw (1,3) node[anchor=east]{A5} -- (3,3);
% Puente de galga
\draw (4,1) node[op amp] (opamp) {};
\draw (opamp.out) -- (5,1) node[anchor=west]{A1};
\end{circuitikz}
\caption{Diagrama simplificado. El puente de galgas requiere alimentación 5V y su salida diferencial va al amplificador.}
\end{figure}

\section*{Preguntas de análisis}
\begin{enumerate}
    \item Compare la sensibilidad del ACS712 (efecto Hall lineal) vs el QMC5883L (magnetómetro triaxial). ¿Cuándo usar cada uno?
    \item ¿Cómo afecta la temperatura a la resistencia de la galga? ¿Cómo se compensa?
    \item Explique la ventaja de usar una PCB sobre protoboard en este tipo de sensores (ruido, impedancia).
    \item ¿Qué fenómeno magnético permite medir corriente sin contacto con el ACS712?
\end{enumerate}

\section*{Bibliografía}
\begin{itemize}
    \item ACS712 Datasheet. Allegro MicroSystems.
    \item QMC5883L Datasheet. QST Corporation.
    \item Galgas extensiométricas – Notas de aplicación Vishay.
\end{itemize}

% !TEX root = 00_instructivo.tex
\chapter{Transductores mecánicos: acelerómetro MPU6050 y medición de altitud}

\section*{Objetivos}
\begin{itemize}
    \item Implementar un contador de pasos usando el acelerómetro MPU6050.
    \item Estimar cambios de altitud mediante doble integración de la aceleración vertical.
    \item Evaluar la precisión del sistema en caminatas reales.
\end{itemize}

\section*{Materiales y equipo}
\begin{itemize}
    \item Módulo MPU6050 (acelerómetro + giroscopio 6 DOF).
    \item Arduino Uno (o Nano) con batería portátil.
    \item Mini protoboard de 9 canales.
    \item Correa para fijar el módulo al tobillo.
    \item Cinta métrica.
\end{itemize}

\section*{Instrucciones}
\begin{enumerate}
    \item Conecte MPU6050: VCC a 5V, GND a GND, SDA a A4, SCL a A5.
    \item Calibre el acelerómetro leyendo los valores en reposo (promedio de 100 muestras).
    \item Sujete el módulo al tobillo. Programe la detección de picos en el eje Z (vertical) con un umbral dinámico.
    \item Realice 5 pruebas de 20 pasos cada una, comparando conteo manual vs conteo del sistema.
    \item Para altitud: integre dos veces la aceleración vertical (después de restar la gravedad). Compare con la altura real medida en escaleras.
    \item Grafique la aceleración en función del tiempo e identifique los picos de paso.
\end{enumerate}

\section*{Esquemático}
\begin{figure}[htbp]
\centering
\begin{circuitikz}[scale=0.8]
\draw (0,0) node[above]{Arduino Uno};
\draw (1,0.5) node[anchor=east]{5V} -- (3,0.5) node[anchor=west]{MPU6050 VCC};
\draw (1,1) node[anchor=east]{GND} -- (3,1) node[anchor=west]{GND};
\draw (1,1.5) node[anchor=east]{A4 (SDA)} -- (3,1.5);
\draw (1,2) node[anchor=east]{A5 (SCL)} -- (3,2);
\draw (4,1.75) node {I2C bus};
\end{circuitikz}
\caption{El MPU6050 se conecta directamente por I2C. Use un capacitor de 0.1µF entre 5V y GND cerca del módulo.}
\end{figure}

\section*{Preguntas de análisis}
\begin{enumerate}
    \item ¿Por qué la doble integración acumula error rápidamente? ¿Cómo se puede mitigar?
    \item ¿Qué papel juega el giroscopio en la corrección de la orientación del acelerómetro?
    \item ¿Cuál es la frecuencia de muestreo mínima recomendada para contar pasos humanos? Justifique.
    \item Proponga un filtro digital (ej. media móvil) para eliminar ruido de alta frecuencia en la señal de aceleración.
\end{enumerate}

\section*{Bibliografía}
\begin{itemize}
    \item MPU6050 Datasheet. InvenSense.
    \item “Step counting using accelerometers” – Application Note, Analog Devices.
    \item Filtro complementario para IMU – Tutorial de Arduino.
\end{itemize}

% !TEX root = 00_instructivo.tex
\chapter{Transductores acústicos: hidrófono con piezoeléctrico y amplificador MAX4466}

\section*{Objetivos}
\begin{itemize}
    \item Construir un hidrófono usando un sensor piezoeléctrico impermeabilizado.
    \item Acondicionar la señal con el amplificador MAX4466.
    \item Clasificar sonidos submarinos (golpes, motores, burbujas) mediante umbrales de amplitud.
\end{itemize}

\section*{Materiales y equipo}
\begin{itemize}
    \item Sensor piezoeléctrico analógico Keyestudio (o disco piezo con cables).
    \item Módulo amplificador MAX4466 (ganancia ajustable).
    \item Arduino Uno.
    \item Mini protoboard de 9 canales.
    \item Recipiente con agua (pecera).
    \item Fuentes de sonido: motor pequeño, objeto metálico, manguera para burbujas.
    \item Silicona o resina epóxica para impermeabilizar.
\end{itemize}

\section*{Instrucciones}
\begin{enumerate}
    \item Cubra el piezoeléctrico con una capa delgada de silicona, dejando el conector expuesto. Deje secar.
    \item Conecte el piezo a la entrada del MAX4466 (IN y GND). La salida del MAX4466 va a un pin analógico del Arduino (A0).
    \item Ajuste la ganancia del MAX4466 para que el ruido de fondo en agua tranquila sea ~1.5V DC (offset).
    \item Genere sonidos: golpee la pared del recipiente, active el motor sumergido, produzca burbujas.
    \item Programe umbrales: si la amplitud pico supera un valor V1 encienda LED rojo, si está entre V0 y V1 LED azul.
    \item Registre 10 segundos de señal para cada fuente y calcule el nivel RMS relativo.
\end{enumerate}

\section*{Esquemático}
\begin{figure}[htbp]
\centering
\begin{circuitikz}[scale=0.7]
\draw (0,0) node[above]{Arduino};
% MAX4466
\draw (2,0) node[op amp, scale=0.8] (max) {};
\draw (max.out) -- (3,0) node[anchor=west]{A0};
\draw (1,-0.5) node[anchor=east]{5V} -- (max.V+);
\draw (1,-1) node[anchor=east]{GND} -- (max.V-);
% Piezo conectado a entrada del MAX4466
\draw (0,0.5) node[pz] (piezo) {};
\draw (piezo.+) -- (max.-);
\draw (piezo.-) -- (max.+);
\end{circuitikz}
\caption{El piezo se conecta a la entrada diferencial del MAX4466. La salida va a A0. Alimentación 5V y GND.}
\end{figure}

\section*{Preguntas de análisis}
\begin{enumerate}
    \item ¿Por qué se necesita un amplificador de ganancia ajustable? ¿Qué pasa si la ganancia es muy alta?
    \item ¿Cómo afecta la impermeabilización a la sensibilidad del piezoeléctrico?
    \item Compare la respuesta en frecuencia de un hidrófono piezoeléctrico vs un micrófono convencional.
    \item Proponga un método para clasificar sonidos no solo por amplitud sino por contenido espectral (FFT).
\end{enumerate}

\section*{Bibliografía}
\begin{itemize}
    \item MAX4466 Datasheet. Maxim Integrated.
    \item “Piezoelectric hydrophone” – Application note, Measurement Specialties.
    \item Keyestudio Analog Piezoelectric Ceramic Vibration Sensor – User manual.
\end{itemize}

\input{06_contaminacion}


\printbibliography

\end{document}
